
\needspace{4\baselineskip}
\color{uniblue}{\subsection{Сборка сигналов (СС)}}
\color{black}

\begin{enumerate}[label=\arabic{section}.\arabic{subsection}.\arabic*, labelsep=4pt, leftmargin=0pt, itemindent=57pt]
%%%%%%%%%%%%%%%%%%%%%%%%%%%%%%%%%%%%%%%%%%%%%%%%%%%%%%%%%%%%%%%%%%%%%%%%% СБОРКА СИГНАЛОВ %%%%%%%%%%%%%%%%%%%%%%%%%%%%%%%%%%%%%%%%%%%%%%%%%%%%%%%%%%%%%%%%%%%%%%%%%%%%%%
\item
В устройстве модуль сборки сигналов представлен функцией «СС». Функция используется в качестве дополнительной логики для предупредительной сигнализации, которая описана в подразделе \ref{ps:punkt}.

\item
Функция принимает различные сигналы от других функций и формирует следующие обобщенные сигналы:
\begin{itemize}
	\item общий сигнал нарушения положения переключателей в выходных цепях «СС: Вых.цепи разобраны»;
	\item общий сигнал нарушения положения испытательных блоков «СС:БИ выведены»;
	\item общий сигнал ошибочного положения коммутационных аппаратов «СС: Неисправость КА»;
	\item общий сигнал превышения времени переключения коммутационных аппаратов «СС: Прев. вр. пер. КА»;
	\item общий сигнал срабатывания внешних функций «СС:Общ. внеш. сигнал».
\end{itemize}

\item
Эти сигналы могут быть использованы для дальнейшего конфигурирования (например, на выходные реле) и применяться в цепях участковых и индивидуальных табло, сигнальных ламп, сигнальных реле, звуковой сигнализации и пр.

\item
Ключом «Вывод терминала» блокируется формирование выходных сигналов.

\item
Логика работы функции выполнена в соответствии с рисунком \ref{ss:img1}. В таблице \ref{ss:tbl1} приведены параметры, необходимые для настройки функции. 

\vspace{2mm}
\begin{figure}[H]
\centering
\includegraphics[width=0.8\textwidth,height=0.8\textheight,keepaspectratio]{\fbpath/216. СС ТН/_latex/img/ss.pdf}
\captionof{figure}{Функциональная схема логики сборки сигналов}\label{ss:img1}
\end{figure}

\end{enumerate}